\documentclass[10pt]{article}

\usepackage[letterpaper,margin=0.66in]{geometry}
\usepackage{amsmath,amssymb}
\usepackage{enumitem}
\usepackage{titlesec}
\usepackage{hyperref}

\hypersetup{colorlinks=true,linkcolor=black,urlcolor=black,citecolor=black}
\titleformat{\section}{\normalsize\bfseries}{}{0em}{}
\titlespacing*{\section}{0pt}{6pt}{2pt}
\setlist[itemize]{nosep,leftmargin=1.2em,topsep=2pt}
\setlength{\parskip}{2.5pt}
\setlength{\parindent}{0pt}
\pagestyle{empty}

\begin{document}

\begin{center}
{\large\bfseries Addendum to ``From the Warp/Wormhole Interface to a Throat-Supported Shift Rail'':\\[2pt]
Geometry Clarification and Current State}\\[5pt]
Daniel S. Goldman \quad Independent researcher \quad
\href{https://orcid.org/0000-0003-2835-3521}{ORCID: 0000-0003-2835-3521}\\[1pt]
September 24, 2026
\end{center}

\section{Clarification}

The reduced diagnostic model of the May 2026 paper \cite{rail} uses the spherically symmetric form
\[
ds^2=-\alpha^2ds^2+\gamma_{ll}\,(dl+\beta^l ds)^2+\gamma_{\Omega\Omega}\,d\Omega^2 ,
\]
with the rail coordinate $l$ as its radial direction. Its implementation sets the areal radius by the Ellis--Morris--Thorne profile
$R^2\propto l^2+R_{th}^2$, a throat of radius $R_{th}=1.75$ at $l=0$. A later revision held $R=1.75$ along the whole
track and opened it into flat space beyond both ends. Either profile joins two asymptotically flat regions through a
minimal sphere. The tested object is therefore a traversable wormhole with a shift service along its throat, in keeping
with the paper's own framing: the wormhole side supplies the plant and the warp side supplies the transport handle.

Three readings in the paper and in the accompanying technical disclosure belong to a rail through one space. The
disclosure's service-time comparison with an exterior null signal takes the rail coordinate as the distance between the
endpoints. The two sides of the throat, however, lie in separate asymptotic regions, so its service-time ratios compare
throat traversal with a reference placed by assumption. Next, the throat relaxation in the service sequence acts on the
support stretch, the support lapse and an angular capacity jacket, while the minimal sphere and the two ends persist.
Finally, the premise of supplied sources holds an existing throat open, whereas forming a two-ended space from ordinary
space requires a change of spatial topology \cite{geroch}.

\section{What carries forward}

The paper's two service rules hold unchanged: transport shift stays inside its support envelope, and the packet is caught
before release. Its service sequence, prepare $\rightarrow$ carry $\rightarrow$ catch $\rightarrow$ release
$\rightarrow$ relax $\rightarrow$ reset, now runs as a literal schedule.

\section{Current state}

The rail now runs along an axis in one flat space \cite{repo}. With unit stretch everywhere,
\[
ds^2=-\alpha^2d\sigma^2+(dz+\beta\,d\sigma)^2+dr^2+r^2d\phi^2,
\qquad
\rho=-\frac{1}{32\pi}\left(\frac{\partial_r\beta}{\alpha}\right)^2 ,
\]
where the energy density follows from the Hamiltonian constraint on the flat spatial slices. The lapse and shift take
their service values within $r\le1.75$ and return to Minkowski values across a transverse boundary layer that ends at
$r=13.25$. The shift carries the packet along a prescribed path with a $2.1c$ lane. A convex lapse plateau of $e^{4}$
around the packet contains the shift, and the live windows open with the lapse leading and lagging the shift.
Meanwhile, a lapse sheath of $e^{1}$ supports the shift's radial transition and travels with the packet on a schedule.
\begin{itemize}
\item \textbf{Arrival.} The packet reaches the track end at $\sigma=3.34$. Light through flat space from the same
entry event arrives at $\sigma=5.0$, which gives the packet a mean speed of $1.35c$.
\item \textbf{Packet.} Its worldtube is timelike from approach to arrival. In the lane the packet moves at $0.016c$
or less relative to local normal observers, and its clock runs at most 55 times exterior time.
\item \textbf{Admissibility.} Over $\sigma\in[-8,10]$ and $|z|\le10$, every non-vacuum boundary-layer point is
Hawking--Ellis Type I, including under refinement, with no Type IV band at any resolved crossing \cite{le}. The
exterior is exactly vacuum.
\item \textbf{Demand.} The standing geometry is flat space. During a transit the demand is a tension without energy on
the shoulders of the lapse profile, local to the packet. Its peak stress is 2.9 and its peak negative energy
$1.9\times10^{-4}$, in units where the service region has radius 1.75. With one unit equal to one metre, the peak
null-energy deficit integrates to about $0.1\,M_\odot$ and the peak stress is $3.5\times10^{44}$\,Pa, both independent
of route length.
\end{itemize}
The next stage evaluates candidate source families against this demand.

\begingroup
\small
\renewcommand{\section}[2]{\vspace{4pt}\textbf{References}\par}
\begin{thebibliography}{9}
\setlength{\itemsep}{0pt}
\bibitem{rail} D.~S. Goldman, ``From the Warp/Wormhole Interface to a Throat-Supported Shift Rail: A Packet-Centered
Active-Rail Architecture after the Wormhole--Warp-Drive Correspondence,'' preprint (2026),
\href{https://doi.org/10.13140/RG.2.2.10402.39368}{doi:10.13140/RG.2.2.10402.39368}.
\bibitem{repo} D.~S. Goldman, active-rail reports, data and code (2026), including \emph{Geometry Clarification: The
Throat-Supported Shift Rail} and \emph{Lapse and Staging Pass},
\url{https://github.com/dgoldman0/spacetime-metric-engineering}.
\bibitem{le} A.~T. Le, ``On the boundary cost of source-consistent warp shells,'' arXiv:2605.25417 (2026).
\bibitem{geroch} R.~P. Geroch, ``Topology in general relativity,'' \emph{J. Math. Phys.} \textbf{8}, 782 (1967).
\end{thebibliography}
\endgroup

\end{document}
